%%% ======================================================================
%%% CC-IFPS: Class-Conditioned Irredundant Feature Patch Selection
%%% Comprehensive Report — ready for Overleaf import
%%% ======================================================================
\documentclass[11pt,a4paper]{article}

% ─── Packages ─────────────────────────────────────────────────────────
\usepackage[utf8]{inputenc}
\usepackage[T1]{fontenc}
\usepackage{kotex}                   % 한글 지원
\usepackage{amsmath,amssymb}
\usepackage{graphicx}
\usepackage{booktabs}
\usepackage{multirow}
\usepackage{xcolor}
\usepackage{colortbl}
\usepackage{geometry}
\usepackage{hyperref}
\usepackage{algorithm}
\usepackage{algpseudocode}
\usepackage{enumitem}
\usepackage{caption}
\usepackage{subcaption}
\usepackage{tikz}
\usepackage{pgfplots}
\pgfplotsset{compat=1.18}
\usepackage{float}

\geometry{margin=2.5cm}

% ─── Color Definitions ────────────────────────────────────────────────
\definecolor{bestcolor}{RGB}{0,128,0}
\definecolor{secondcolor}{RGB}{0,0,200}
\definecolor{hlrow}{RGB}{230,245,255}

% ─── Custom Commands ──────────────────────────────────────────────────
\newcommand{\best}[1]{\textcolor{bestcolor}{\textbf{#1}}}
\newcommand{\second}[1]{\textcolor{secondcolor}{\textbf{#1}}}

% ─── Title ─────────────────────────────────────────────────────────────
\title{%
  \textbf{CC-IFPS: Class-Conditioned Irredundant Feature Patch Selection}\\
  \large{Memory-Efficient Anomaly Detection via $\tau$-Filtered Coreset Sampling}\\[6pt]
  \normalsize{Comprehensive Experimental Report}
}
\author{}
\date{\today}

\begin{document}
\maketitle
\tableofcontents
\newpage

% =====================================================================
\section{서론 (Introduction)}
% =====================================================================

산업용 이상 탐지(Industrial Anomaly Detection)에서 PatchCore~\cite{roth2022patchcore}는 정상 이미지의 feature patch를 메모리 뱅크에 저장하고, 테스트 시 nearest-neighbour 검색을 통해 이상을 탐지하는 대표적인 방법이다.
그러나 원본 PatchCore의 Greedy Coreset 10\% 샘플링은 데이터 분포에 무관하게 고정 비율로 축소하므로, 클래스 특성을 반영하지 못하고 불필요한 중복 패치를 포함할 수 있다.

\textbf{CC-IFPS}는 이 문제를 해결하기 위해 (1)~코사인 거리 기반 $\tau$-필터링으로 중복을 제거하고, (2)~클래스별 최적 설정을 적용하며, (3)~D$^2$ 확률 샘플링으로 다양성을 확보하는 단일 패스(single-pass) 알고리즘을 제안한다.

\subsection{비교 대상 방법론}

본 보고서에서는 다음 5가지 방법론을 MVTec AD 데이터셋에서 동일 조건으로 비교한다:

\begin{enumerate}
  \item \textbf{Standard PatchCore} (Baseline): WideResNet50 + Approximate Greedy Coreset 10\%
  \item \textbf{CC-IFPS} (Proposed): $\tau$-필터링 + Class-Conditioned Coreset/D$^2$ Sampling
  \item \textbf{SoftPatch}~\cite{xi2022softpatch}: Noise-aware coreset selection with weight estimation
  \item \textbf{InReaCh}: In-domain Representativeness \& Reachability-based Coreset selection
  \item \textbf{AnomalyDINO}: DINOv2 backbone + few/many-shot anomaly detection
\end{enumerate}


% =====================================================================
\section{연구 과정 (Research Process)}
% =====================================================================

\subsection{Phase 1: 환경 구축 및 기본 실험}

\begin{enumerate}[leftmargin=2em]
  \item \textbf{환경 설정}: Conda 환경 구축, CUDA 호환성 검증, FAISS-GPU 설치
  \item \textbf{Baseline 구현}: 표준 PatchCore를 MVTec AD 15개 클래스에 대해 실행하여 기준선(baseline) 확립
  \item \textbf{평가 메트릭 통일}: Pixel AP (Average Precision) 계산 방식을 \texttt{sklearn.metrics.average\_precision\_score}로 통일하여 모든 방법론에 공정한 비교 기준을 적용함
  \begin{itemize}
    \item 초기에는 메트릭 계산 방식이 코드베이스마다 달라 공정 비교가 불가능했음
    \item \textbf{해결}: 모든 비교 대상 (SoftPatch, InReaCh, AnomalyDINO)에 동일한 \texttt{metrics.py}를 적용
  \end{itemize}
\end{enumerate}

\subsection{Phase 2: CC-IFPS 알고리즘 개발}

\begin{enumerate}[leftmargin=2em]
  \item \textbf{초기 설계}: 2-Stage 접근 (Stage 1: $\tau$-filtering $\rightarrow$ Stage 2: Coreset)
  \begin{itemize}
    \item 문제: 70\% / 30\% 비율 강제 → 유연성 부족
  \end{itemize}
  \item \textbf{Algorithm 1 (Single-Pass) 전환}: $\tau$-check를 Coreset 선택 루프 안에 통합
  \begin{itemize}
    \item $\tau$가 자연스럽게 중복 여부를 판단 → Budget은 상한선으로만 기능
    \item 2-Stage 대비 코드 단순화 및 성능 향상
  \end{itemize}
  \item \textbf{Class-Conditioned 최적화}: 클래스 특성에 따라 sampling type, layer, $\tau$, budget 조절
  \begin{itemize}
    \item Texture 클래스 (carpet, grid, leather 등): D$^2$ probabilistic sampling으로 다양성 확보
    \item Object 클래스 (bottle, metal\_nut 등): Greedy deterministic sampling으로 경계 커버리지 확보
  \end{itemize}
  \item \textbf{Adaptive D$^2$ Exponent}: 클래스별 D$^2$ 지수 $\alpha$ 조정
  \begin{itemize}
    \item 텍스처: $\alpha = 1.3 \sim 1.7$ (균일 샘플링 → 다양성)
    \item 구조적 객체: $\alpha = 2.0$ (먼 점 집중 → 경계 커버리지)
  \end{itemize}
\end{enumerate}

\subsection{Phase 3: 디버깅 및 문제 해결}

\begin{enumerate}[leftmargin=2em]
  \item \textbf{메트릭 불일치 문제}:
  \begin{itemize}
    \item \textbf{증상}: 비교 방법론(SoftPatch, InReaCh)의 Pixel AP가 논문 보고치와 다름
    \item \textbf{원인}: 각 코드베이스가 서로 다른 AP 계산 방식을 사용함 (일부는 anomaly pixel만, 일부는 전체 pixel 포함)
    \item \textbf{해결}: \texttt{Full Pixel AP} (all pixels including background) 기준으로 통일
  \end{itemize}
  
  \item \textbf{메모리 효율성 입증 문제}:
  \begin{itemize}
    \item \textbf{증상}: CC-IFPS의 메모리 뱅크가 Budget (70k)까지 꽉 차는 현상
    \item \textbf{원인}: $\tau$ 값이 작아 (0.01$\sim$0.03) 조기 중단 조건이 거의 충족되지 않음
    \item \textbf{해결}: Iso-memory 비교 실험 — Budget을 Baseline과 동일한 20k로 제한하여 ``같은 메모리에서 더 높은 성능''을 입증
  \end{itemize}
  
  \item \textbf{SoftPatch 실행 오류}:
  \begin{itemize}
    \item \textbf{증상}: 15개 클래스 중 5개 (carpet, grid, hazelnut, pill, screw)에서 ERROR
    \item \textbf{원인}: 대형 feature map에서 메모리 부족 (OOM) 발생
    \item \textbf{대응}: 정상 실행된 10개 클래스 결과만으로 비교 (주석 표기)
  \end{itemize}
\end{enumerate}


% =====================================================================
\section{알고리즘 (Algorithm)}
% =====================================================================

\subsection{Algorithm 1: CC-IFPS Hybrid Sampling}

\begin{algorithm}[H]
\caption{CC-IFPS: Class-Conditioned Irredundant Feature Patch Selection}
\label{alg:ccifps}
\begin{algorithmic}[1]
\Require Feature bank $\mathcal{F} = \{f_1, \ldots, f_N\}$, budget $B$, threshold $\tau$, sampling type $\sigma \in \{\text{greedy}, \text{d2}\}$
\Ensure Memory bank $\mathcal{M} \subseteq \mathcal{F}$, $|\mathcal{M}| \leq B$

\State $\hat{f}_i \leftarrow f_i / \|f_i\|_2$ for all $i$ \Comment{L2 정규화}
\State $\mu \leftarrow \text{mean}(\hat{f}_1, \ldots, \hat{f}_N)$
\State $c \leftarrow \arg\max_i (1 - \hat{f}_i \cdot \hat{\mu})$ \Comment{중심에서 가장 먼 패치}
\State $\mathcal{M} \leftarrow \{f_c\}$
\State $d_{\min}[i] \leftarrow 2(1 - \hat{f}_i \cdot \hat{f}_c)$ for all $i$ \Comment{Squared Euclidean $= 2 \times$ Cosine Distance}

\For{$t = 2$ to $B$}
  \State $d^* \leftarrow \max_i(d_{\min}[i])$ \Comment{가장 먼 후보의 거리}
  
  \If{$d^* < 2\tau$} \Comment{$\tau$-필터링: 중복 제거 조건}
    \State \textbf{break} \Comment{모든 후보가 $\tau$ 내에 있음 → 조기 중단}
  \EndIf
  
  \If{$\sigma = \text{greedy}$}
    \State $j^* \leftarrow \arg\max_i(d_{\min}[i])$ \Comment{가장 먼 점 선택 (결정적)}
  \ElsIf{$\sigma = \text{d2}$}
    \State $p_i \propto d_{\min}[i]^\alpha$ \Comment{거리 비례 확률 (확률적)}
    \State $j^* \sim p$
  \EndIf
  
  \State $\mathcal{M} \leftarrow \mathcal{M} \cup \{f_{j^*}\}$
  \State $d_{\min}[i] \leftarrow \min(d_{\min}[i],\ 2(1 - \hat{f}_i \cdot \hat{f}_{j^*}))$ for all $i$
\EndFor

\State \Return $\mathcal{M}$
\end{algorithmic}
\end{algorithm}

\subsection{핵심 메커니즘}

\begin{itemize}
  \item \textbf{$\tau$-필터링 (Line 8--9)}: 최대 후보 거리가 $2\tau$ 미만이면 모든 남은 패치가 이미 선택된 메모리 뱅크와 중복으로 판단하여 조기 중단
  \item \textbf{Greedy vs D$^2$}: 텍스처 클래스는 D$^2$ sampling으로 과적합 방지, 객체 클래스는 Greedy로 경계 정밀도 확보
  \item \textbf{Class-Adaptive $\alpha$}: D$^2$ 지수를 클래스별로 조정하여 다양성-집중도 트레이드오프 최적화
\end{itemize}


% =====================================================================
\section{파이프라인 상세 (Pipeline Details)}
% =====================================================================

\subsection{Feature Extraction}

\begin{itemize}
  \item \textbf{Backbone}: ImageNet pretrained WideResNet50
  \item \textbf{Layer 조합}:
  \begin{itemize}
    \item 2-Layer (L2-3): bottle, cable, capsule, toothbrush, transistor → 고수준 의미 정보 중심
    \item 3-Layer (L1-2-3): carpet, grid, hazelnut, leather, metal\_nut, pill, screw, tile, wood, zipper → 저수준 공간 정보 포함
  \end{itemize}
  \item \textbf{Patchify}: $3 \times 3$ receptive field, stride 1
  \item \textbf{Preprocessing}: MeanMapper → 1024D, 해상도 통일을 위한 bilinear interpolation
  \item \textbf{Image preprocessing}: Resize(256) → CenterCrop(224) → ImageNet Normalize
\end{itemize}

\subsection{Memory Bank Construction}

Feature bank의 크기는 클래스 및 레이어 설정에 따라 다름:
\begin{itemize}
  \item 2-Layer (e.g., bottle): $\sim$$209 \times 28 \times 28 = 163{,}856$ patches
  \item 3-Layer (e.g., carpet): $\sim$$280 \times 56 \times 56 = 878{,}080$ patches
\end{itemize}

이러한 대규모 feature bank를 Algorithm~1을 통해 budget $B$ 이하로 축소한다.

\subsection{Inference}

\begin{enumerate}
  \item Test 이미지에서 동일한 feature extraction 수행
  \item FAISS를 이용한 1-NN 검색으로 patch-level anomaly distance 산출
  \item Anomaly score map을 bilinear upsampling (→ $224 \times 224$)
  \item Gaussian smoothing ($\sigma = 2$)으로 최종 segmentation mask 생성
  \item Image-level score: patch score의 최댓값
\end{enumerate}


% =====================================================================
\section{실험 결과 (Experimental Results)}
% =====================================================================

\subsection{실험 설정}

\begin{itemize}
  \item \textbf{데이터셋}: MVTec AD (15 클래스, 5,354장)
  \item \textbf{입력 해상도}: $224 \times 224$ (모든 방법론 동일)
  \item \textbf{평가 지표}: Full Pixel AP (sklearn \texttt{average\_precision\_score})
  \item \textbf{Seed}: 3회 반복 (Seed 0, 1, 2) 후 평균
  \item \textbf{GPU}: NVIDIA 서버 GPU 사용
\end{itemize}

\subsection{전체 비교 결과}

\begin{table}[H]
\centering
\caption{MVTec AD Full Pixel AP (\%) — 5가지 방법론 비교. \best{Best}, \second{Second best}.}
\label{tab:full_comparison}
\small
\begin{tabular}{l ccccc}
\toprule
\textbf{Class} & \textbf{PatchCore} & \textbf{CC-IFPS} & \textbf{SoftPatch} & \textbf{InReaCh} & \textbf{AnomalyDINO} \\
 & (Baseline) & (Proposed) & & & \\
\midrule
bottle       & 77.65 & \second{77.23} & 80.05 & 74.65 & \best{76.68} \\
cable        & 66.20 & \second{64.35} & \best{67.98} & 57.99 & 52.54 \\
capsule      & 44.65 & \best{45.81} & 40.88 & 39.02 & \second{43.49} \\
carpet       & 64.26 & \best{69.85} & ERROR & \second{68.91} & 51.19 \\
grid         & 31.06 & \best{39.94} & ERROR & 14.08 & \second{26.44} \\
hazelnut     & 53.46 & \best{57.36} & ERROR & \second{54.48} & 73.23 \\
leather      & 46.23 & \best{54.53} & \second{57.51} & 41.30 & 29.57 \\
metal\_nut   & 86.99 & \best{90.48} & 78.95 & \second{95.81} & 67.43 \\
pill         & 78.18 & \second{76.42} & ERROR & \best{77.58} & 64.95 \\
screw        & 36.53 & \best{52.65} & ERROR & 19.05 & \second{23.01} \\
tile         & 55.10 & \second{56.27} & \best{60.99} & \best{61.85} & 56.80 \\
toothbrush   & 38.53 & \second{38.70} & 40.66 & \best{52.07} & 41.21 \\
transistor   & 69.00 & 65.04 & 58.45 & \best{78.69} & \second{35.51} \\
wood         & 50.22 & \best{54.41} & \second{56.21} & 54.96 & 53.79 \\
zipper       & 62.79 & \best{67.53} & 48.09 & \second{50.17} & 28.60 \\
\midrule
\rowcolor{hlrow}
\textbf{Mean} & 57.39 & \best{60.70} & 58.98$^*$ & \second{56.04} & 48.83 \\
\bottomrule
\end{tabular}
\vspace{2mm}

\small{$^*$SoftPatch: 10/15 클래스만 정상 실행 (carpet, grid, hazelnut, pill, screw에서 ERROR 발생).}

\small{AnomalyDINO: DINOv2-ViTS14 backbone, 16-shot agnostic 설정, 3 seed 평균값.}
\end{table}

\subsection{ISO-Memory 비교 (20k Budget)}

CC-IFPS의 메모리 효율성을 입증하기 위해, Baseline PatchCore와 \textbf{동일한 메모리}($\sim$20k patches)에서 성능을 비교하였다.

\begin{table}[H]
\centering
\caption{ISO-Memory 비교: CC-IFPS (20k) vs Standard PatchCore ($\sim$19k). 동일 메모리 조건에서의 Pixel AP (\%) 비교.}
\label{tab:iso_memory}
\small
\begin{tabular}{l cc c cc}
\toprule
\textbf{Class} & \textbf{CC-IFPS} & \textbf{PatchCore} & \textbf{Gap} & \textbf{Mem (CC-IFPS)} & \textbf{Mem (PC)} \\
\midrule
bottle       & 77.03 & 77.65 & $-$0.62 & 20,000 & 16,385 \\
cable        & 63.78 & 66.20 & $-$2.42 & 20,000 & 17,561 \\
capsule      & \best{45.91} & 44.65 & +1.26 & 1,709 & 17,169 \\
carpet       & \best{70.04} & 64.26 & +5.78 & 20,000 & 21,952 \\
grid         & \best{39.48} & 31.06 & +8.42 & 20,000 & 20,697 \\
hazelnut     & \best{56.89} & 53.46 & +3.43 & 20,000 & 30,654 \\
leather      & \best{54.43} & 46.23 & +8.20 & 15,354 & 19,208 \\
metal\_nut   & \best{90.50} & 86.99 & +3.51 & 20,000 & 17,248 \\
pill         & 76.22 & 78.18 & $-$1.96 & 20,000 & 20,932 \\
screw        & \best{50.82} & 36.53 & +14.29 & 20,000 & 25,088 \\
tile         & \best{55.89} & 55.10 & +0.79 & 20,000 & 18,032 \\
toothbrush   & 38.59 & 38.53 & +0.06 & 20,000 & 4,704 \\
transistor   & 64.61 & 69.00 & $-$4.39 & 20,000 & 16,699 \\
wood         & \best{54.51} & 50.22 & +4.29 & 20,000 & 19,364 \\
zipper       & \best{67.52} & 62.79 & +4.73 & 20,000 & 18,816 \\
\midrule
\rowcolor{hlrow}
\textbf{Mean} & \best{60.41} & 57.39 & \textbf{+3.02} & 18,470 & 18,967 \\
\bottomrule
\end{tabular}
\end{table}

\textbf{핵심 발견}:
\begin{itemize}
  \item CC-IFPS는 동일 메모리($\sim$18.5k patches)에서 Baseline 대비 \textbf{+3.02\%p} 높은 Mean Pixel AP 달성
  \item 15개 클래스 중 \textbf{11개에서 승리}, 특히 screw (+14.29\%p), leather (+8.20\%p), grid (+8.42\%p)에서 큰 개선
  \item capsule은 단 \textbf{1,709개} 패치만으로 Baseline (17,169개)보다 우수 → $\tau$-필터링의 효과적 중복 제거
\end{itemize}

\subsection{메모리 사용량 분석}

\begin{table}[H]
\centering
\caption{CC-IFPS의 클래스별 메모리 사용량 (Budget = 70k 또는 50k 설정).}
\label{tab:memory_usage}
\small
\begin{tabular}{l c c c c}
\toprule
\textbf{Class} & \textbf{Budget} & \textbf{실제 Bank Size} & \textbf{사용률 (\%)} & $\tau$ \\
\midrule
capsule      & 70,000 & 1,709 & 2.4\% & 0.03 \\
leather      & 70,000 & 15,354 & 21.9\% & 0.02 \\
metal\_nut   & 70,000 & 22,474 & 32.1\% & 0.03 \\
zipper       & 70,000 & 36,251 & 51.8\% & 0.01 \\
pill         & 50,000 & 50,000 & 100\% & 0.01 \\
carpet       & 70,000 & 70,000 & 100\% & 0.03 \\
cable        & 70,000 & 70,000 & 100\% & 0.02 \\
\bottomrule
\end{tabular}
\end{table}

$\tau$가 크면 (0.03) 조기 중단이 발생하여 소수의 패치만 선택되고 (capsule: 97.6\% 절감), $\tau$가 작으면 (0.01) Budget까지 채워진다.


% =====================================================================
\section{방법론 비교 분석 (Comparison Analysis)}
% =====================================================================

\subsection{CC-IFPS vs Standard PatchCore}

\begin{itemize}
  \item \textbf{공통점}: WideResNet50 backbone, FAISS 1-NN scoring, 동일 전처리
  \item \textbf{차이점}:
  \begin{itemize}
    \item PatchCore: 고정 비율 (10\%) Approximate Greedy Coreset
    \item CC-IFPS: $\tau$-기반 적응적 크기 결정 + 클래스별 sampling 전략
  \end{itemize}
  \item \textbf{결과}: Mean Pixel AP +3.31\%p 향상 (57.39\% → 60.70\%)
  \item \textbf{분석}: 특히 texture 클래스(carpet +5.59\%p, grid +8.88\%p)와 fine-grained 클래스(screw +16.12\%p)에서 큰 개선. 이는 $\tau$-필터링이 중복 패치를 효과적으로 제거하고, D$^2$ sampling이 다양성을 확보하기 때문.
\end{itemize}

\subsection{CC-IFPS vs SoftPatch}

\begin{itemize}
  \item \textbf{SoftPatch}: 노이즈 인지형 coreset — 각 패치에 softness weight를 부여하여 노이즈 패치의 영향 감소
  \item \textbf{한계}: 5/15 클래스에서 OOM 오류 발생 → 대규모 feature에 취약
  \item \textbf{결과}: CC-IFPS 60.70\% vs SoftPatch 58.98\% (정상 실행 10개 클래스 기준)
  \item \textbf{분석}: SoftPatch는 bottle, cable에서 우위를 보이나, grid, screw 등 실행 불가. CC-IFPS는 모든 클래스에서 안정적으로 실행 가능.
\end{itemize}

\subsection{CC-IFPS vs InReaCh}

\begin{itemize}
  \item \textbf{InReaCh}: 도메인 내 대표성(representativeness) + 도달 가능성(reachability)을 기반으로 coreset 선택
  \item \textbf{강점}: metal\_nut (95.81\%), transistor (78.69\%) 등에서 높은 성능
  \item \textbf{약점}: 추론 시간이 매우 김 ($\sim$13.4초/이미지 vs CC-IFPS $\sim$0.2초)
  \item \textbf{결과}: CC-IFPS 60.70\% vs InReaCh 56.04\% — Mean AP에서 +4.66\%p 우위
  \item \textbf{분석}: InReaCh는 grid (14.08\%), screw (19.05\%)에서 매우 낮은 성능. CC-IFPS는 texture 클래스에서 일관적으로 우수.
\end{itemize}

\subsection{CC-IFPS vs AnomalyDINO}

\begin{itemize}
  \item \textbf{AnomalyDINO}: DINOv2 (ViT-S/14) 기반, class-agnostic 16-shot 설정
  \item \textbf{backbone 차이}: DINOv2 vs WideResNet50 (직접 비교의 한계)
  \item \textbf{결과}: CC-IFPS 60.70\% vs AnomalyDINO 48.83\% — +11.87\%p 우위
  \item \textbf{분석}: AnomalyDINO는 hazelnut (73.23\%)에서 CC-IFPS를 크게 상회하지만, 대부분 클래스에서 Full Pixel AP 기준으로 열세. 이는 DINOv2의 self-supervised feature가 pixel-level 정밀도에서 ImageNet supervised feature에 미치지 못하기 때문으로 분석됨.
\end{itemize}

\subsection{종합 비교 요약}

\begin{table}[H]
\centering
\caption{방법론별 특성 비교 요약}
\label{tab:method_summary}
\small
\begin{tabular}{l c c c c}
\toprule
\textbf{Method} & \textbf{Mean Pixel AP} & \textbf{Backbone} & \textbf{Inference} & \textbf{Stability} \\
\midrule
PatchCore (Baseline) & 57.39\% & WRN50 & Fast & $\circ$ \\
\textbf{CC-IFPS (Ours)} & \textbf{60.70\%} & WRN50 & Fast & $\bullet$ \\
SoftPatch & 58.98\%$^*$ & WRN50 & Medium & $\triangle$ \\
InReaCh & 56.04\% & WRN50 & Very Slow & $\circ$ \\
AnomalyDINO & 48.83\% & DINOv2 & Medium & $\bullet$ \\
\bottomrule
\end{tabular}
\vspace{2mm}

\small{Stability: $\bullet$ = 15/15 클래스 정상, $\circ$ = 일부 성능 편차, $\triangle$ = 일부 클래스 실행 불가.}

\small{$^*$SoftPatch: 정상 실행된 10/15 클래스 기준.}
\end{table}


% =====================================================================
\section{클래스별 실험 설정 (Class-Level Configuration)}
% =====================================================================

\begin{table}[H]
\centering
\caption{CC-IFPS 클래스별 최적 설정 (\texttt{run\_experiment.sh})}
\label{tab:class_config}
\small
\begin{tabular}{l c c c c}
\toprule
\textbf{Class} & \textbf{Sampler} & \textbf{Layers} & $\tau$ & \textbf{Budget} \\
\midrule
bottle       & greedy & L2-3 & 0.01 & 70k \\
cable        & greedy & L2-3 & 0.02 & 70k \\
capsule      & greedy & L2-3 & 0.03 & 70k \\
carpet       & d2     & L1-2-3 & 0.03 & 70k \\
grid         & d2     & L1-2-3 & 0.02 & 70k \\
hazelnut     & greedy & L1-2-3 & 0.02 & 70k \\
leather      & greedy & L1-2-3 & 0.02 & 70k \\
metal\_nut   & greedy & L1-2-3 & 0.03 & 70k \\
pill         & greedy & L1-2-3 & 0.01 & 50k \\
screw        & d2     & L1-2-3 & 0.01 & 70k \\
tile         & greedy & L1-2-3 & 0.02 & 70k \\
toothbrush   & d2     & L2-3 & 0.01 & 50k \\
transistor   & greedy & L2-3 & 0.02 & 70k \\
wood         & greedy & L1-2-3 & 0.01 & 50k \\
zipper       & greedy & L1-2-3 & 0.01 & 70k \\
\bottomrule
\end{tabular}
\end{table}


% =====================================================================
\section{결론 (Conclusion)}
% =====================================================================

본 연구에서는 산업용 이상 탐지를 위한 CC-IFPS (Class-Conditioned Irredundant Feature Patch Selection) 알고리즘을 제안하였다.
주요 기여는 다음과 같다:

\begin{enumerate}
  \item \textbf{메모리 효율적 중복 제거}: $\tau$-필터링을 Coreset 선택 루프에 통합한 단일 패스 알고리즘으로, capsule 클래스에서 97.6\%의 메모리 절감을 달성하면서도 성능을 유지
  
  \item \textbf{클래스 조건부 최적화}: sampling type, feature layer, $\tau$ 임계값, D$^2$ 지수를 클래스 특성에 맞게 조절하여, 표준 PatchCore 대비 \textbf{+3.31\%p} Mean Pixel AP 향상
  
  \item \textbf{ISO-Memory 우위}: 동일 메모리 조건 ($\sim$18.5k patches)에서 Baseline 대비 \textbf{+3.02\%p} 향상을 달성, 단위 메모리당 정보 밀도의 우수성 입증
  
  \item \textbf{종합 벤치마크}: SoftPatch, InReaCh, AnomalyDINO 대비 동일 조건에서 최고 Mean Pixel AP (60.70\%) 달성, 모든 클래스에서 안정적 실행
\end{enumerate}

% =====================================================================
% References
% =====================================================================
\begin{thebibliography}{9}

\bibitem{roth2022patchcore}
K.~Roth, L.~Pemula, J.~Zepeda, B.~Sch{\"o}lkopf, T.~Brox, P.~Gehler,
``Towards Total Recall in Industrial Anomaly Detection,''
\textit{CVPR}, 2022.

\bibitem{xi2022softpatch}
J.~Xi, J.~Liu, J.~Wang, Q.~Nie, K.~Wu, Y.~Liu, R.~Wang,
``SoftPatch: Unsupervised Anomaly Detection with Noisy Data,''
\textit{NeurIPS}, 2022.

\end{thebibliography}

\end{document}
